\documentclass{chsmc}
\chsmcsetup{
	year = 2014,
	part = 1,
	type = A,
	show-answers = false,
}
\begin{document}

\section{填空题：本大题共 8 小题，每小题 8 分，共 64 分.}

\begin{enumerate}
	\item 若正数 $a$, $b$ 满足
		$2 + \log_2 a = 3 + \log_3 b = \log_6(a+b)$.
		则 $\dfrac{1}{a} + \dfrac{1}{b}$ 的值为 \blankline
	\item 设集合
		$\left\{\left.\dfrac{3}{a} + b\,\right|\, 1 \leqslant a \leqslant b \leqslant 2\right\}$
		中的最大元素与最小元素分别为 $M$, $m$，则 $M-m$ 的值为 \blankline
	\item 若函数 $f(x) = x^2 + a|x-1|$ 在 $[0,+\infty)$ 上单调递增，
		则实数 $a$ 的取值范围是 \blankline
	\item 数列 $\{a_n\}$ 满足 $a_1 = 2$，
		$a_{n+1} = \dfrac{2(n+2)}{n+1}a_n$（$n \in \mathbf{N}^*$），则
		$\dfrac{a_{2014}}{a_1 + a_2 + \cdots + a_{2013}} =$ \blankline
	\item 正四棱锥 $P$-$ABCD$ 中，侧面是边长为 $1$ 的正三角形，
		$M$, $N$ 分别是边 $AB$, $BC$ 的中点，则异面直线 $MN$ 与 $PC$ 之间的距离是 \blankline
	\item 设椭圆 $\Gamma$ 的两个焦点是 $F_1$, $F_2$，过点 $F_1$ 的直线与 $\Gamma$ 交于点 $P$, $Q$.
		若 $|PF_1| = |F_1F_2|$，且 $3|PF_1| = 4|QF_1|$，
		则椭圆 $\Gamma$ 的短轴与长轴的比值为 \blankline
	\item 设等边三角形 $ABC$ 的内切圆半径为 $2$，圆心为 $I$.
		若点 $P$ 满足 $PI = 1$，则 $\triangle APB$ 与 $\triangle APC$ 的面积之比的最大值为 \blankline
	\item 设 $A$, $B$, $C$, $D$ 是空间四个不共面的点，以 $\dfrac{1}{2}$ 的概率在每对点之间连一条边，
		任意两对点之间是否连边是相互独立的，则 $A$, $B$ 可用
		（一条边或者若干条边组成的）空间折线连接的概率为 \blankline
\end{enumerate}

\section{解答题：本大题共 3 个小题，满分 56 分. 解答应写出文字说明、证明过程或演算步骤.}

\begin{enumerate}[resume]
	\item (本题满分 16 分) 平面直角坐标系 $xOy$ 中，$P$ 是不在 $x$ 轴上的一个动点，满足条件：
		过 $P$ 可作抛物线 $y^2 = 4x$ 的两条切线，两切点连线 $l_P$ 与 $PO$ 垂直.
		设直线 $l_P$ 与直线 $PO$, $x$ 轴的交点分别为 $Q$, $R$.
		\begin{enumerate}[(1)]
			\item 证明 $R$ 是一个定点;
			\item 求 $\dfrac{|PQ|}{|QR|}$ 的最小值.
		\end{enumerate}
	\item (本题满分 20 分) 数列 $\{a_n\}$ 满足 $a_1 = \dfrac{\pi}{6}$，
		$a_{n+1} = \arctan(\sec a_n)$（$n \in \mathbf{N}^*$）.
		求正整数 $m$，使得 $\sin a_1\cdot\sin a_2\cdot\cdots\cdot\sin a_m = \dfrac{1}{100}$.
	\item (本题满分 20 分) 确定所有的复数 $\alpha$，使得对任意复数 $z_1$, $z_2$
		（$|z_1|$, $|z_2| < 1$，$z_1 \neq z_2$），均有
		$(z_1+\alpha)^2 + \alpha\overline{z_1} \neq (z_2+\alpha)^2 + \alpha\overline{z_2}$.
\end{enumerate}

\end{document}
